\textbf{Trefethen 3.5}

Using the commands \texttt{tic} and \texttt{toc}, study the time taken by \texttt{MATLAB} to compute an $N$-point FFT,
as a function of $N$. If $N = 2^k$ for $k = 0, 1, ...,15$, for example, what is the dependence of the time on $N$? What
if $N = 500, 501, ..., 519, 520$? From a plot of the latter results, can you spot the prime numbers in this range?


\begin{solution}
  For powers of 2, we observe a linear dependence between the time taken to compute the FFT and $N$:

  \begin{figure}[h]
    \centering
    \includegraphics[width=\textwidth]{problem_6.png}
    \caption{FFT Computation Time vs. $N$}
  \end{figure}

  The four most time-consuming values of $N$ occur at $N = 501, 503, 509, 519$. Of these, 501 and 519 are the product of
  a prime number and 3; 503 and 509, on the other hand, are prime numbers. We further observe that the least 
  time-consuming FFTs happen at multiples of two.
  \ \\
\end{solution}